\subsection{HOMEOMORPHISMS OF AN INTERVAL ONTO AN INTERVAL}

THEOREM 5. Let I be an interval in R. Then a mapping f of I into R is a homeomorphism of I onto \(f(I)\) if and only if f is strictly monotonic and continuous on I; and \(f(I)\) is then an interval in R.

\begin{enumerate}
\def\labelenumi{\arabic{enumi})}
\tightlist
\item
  The condition is necessary. Let \(a\) and \(b\) be two points of \(I\) such that \(a < b\) , and suppose for example that \(f(a) < f(b)\) . Let us show that \(f\) is strictly increasing on \(I\) . First, if \(a < c < b\) then we must have \(f(a) < f(c) < f(b)\) ; if, for example, we had \(f(a) < f(b) < f(c)\)
\end{enumerate}

then the image of the interval \([a, c]\) under \(f\) would be a connected set (Chapter I, § 11, no. 2, Proposition 4) and would therefore contain the interval \([f(a), f(c)]\) ; hence there would exist \(x \in [a, c]\) such that \(f(x) = f(b)\) , contrary to the hypothesis that \(f\) is injective.

It follows that if \(x\) and \(y\) are any two points of \(I\) such that \(x < y\) , then \(f(x) < f(y)\) ; for we have \(f(a) < f(x) < f(b)\) if \(a < x < b\) , \(f(a) < f(b) < f|(x)\) if \(b < x\) , and \(f(x) < f(a) < f(b)\) if \(x < a\) ; repeating the argument with \(a, x, y\) in place of \(a, b, x\) respectively, we see that \(f(x) < f(y)\) .

\begin{enumerate}
\def\labelenumi{\arabic{enumi})}
\setcounter{enumi}{1}
\tightlist
\item
  The condition is sufficient. Suppose that \(f\) is continuous and strictly monotonic on \(I\) (say, strictly increasing): \(f(I)\) is connected and is therefore an interval, and since \(f\) is strictly increasing, \(f\) is a bijective mapping of \(I\) onto \(f(I)\) . Moreover, the image under \(f\) of an open interval in \(I\) is an open interval in \(f(I)\) , and therefore (§1, no, 4, Proposition 5) \(f\) is a homeomorphism of \(I\) onto \(f(I)\) .
\end{enumerate}

Remark. The first part of the preceding proof shows in fact that a continuous injective mapping of I into R is strictly monotonic; from the second part of the proof it therefore follows that every continuous injective mapping f of an interval I into R is a homeomorphism of I onto \(f(\mathbf{I})\) .

